\section{Discussion and future directions \label{sec:discussion}}

The results of Sec.~\ref{sec:results} demonstrate clearly that the light
hadrons can be studied using operators that excite only a very small number of
extended degrees of freedom on each time-slice. The restricted set of modes
needed to create a hadron means all elements of quark propagation needed to
construct the correlation function can be measured directly. This does not
preclude using stochastic methods \cite{Foley:2005ac} to reduce computational
costs.

In the straightforward recipe described in this paper, the distillation space
was constructed from the lowest eigenvectors of the lattice three-dimensional
Laplace operator on a time-slice. This definition is not unique and more
experimentation with different choices of distillation space might prove
useful. Also, the operator used in this work is diagonal in the quark spin
index. Alternative choices of the distillation space might include non-trivial
spin structure, using for example the vector space spanned by the lowest modes
of the three-dimensional Dirac operator.

Recall that in the simplest implementation tested here when $N=M$,
the distillation operator becomes the identity and no smearing is performed.
This is clearly an undesirable feature that is circumvented by the obvious
variation on the method to include different weights for the eigenmodes. This
more general distillation operator is then
\begin{equation}
        {\cal J}(t) = \sum_{k=1}^N v^{(k)}(t) \, f\big(\lambda_k(t)\big)
              \, v^{(k)\dag}(t).
\end{equation}
Including weight functions on the eigenvectors can be used to suppress
more rapidly fluctuating contributions when many eigenvectors are employed.
Using exponential weights leads to increasingly accurate approximations to the
conventional smearing algorithm of Eq.~\ref{eqn:smearing}.

No increase in the number of eigenvectors would be expected as the continuum
limit is taken at a fixed physical volume since, as is seen in the data of
Sec.~\ref{subsec:results-wavefn}, the smooth distilled modes are
almost completely decoupled from the cut-off dynamics. This expectation is yet
to be tested numerically.
One substantial possible drawback with the method is the need to increase the
rank of the distillation operator as the three-dimensional lattice volume
increases. The corresponding linear increase in the cost of solving the Dirac
equation leads to an algorithm for spectroscopy that scales like ${\cal
O}(V^2)$.
It appears from the results we have presented that with the required
increase in rank with increasing volume one also obtains a decrease in
statistical fluctuations.  The
possible solution to the issue of volume scaling may be to combine the method
with a suitably designed stochastic estimation scheme. This is under
investigation.

Note that the method does not give direct access to all elements of the quark
propagator, only those relevant to low-energy spectroscopy. This would be a
limitation in calculations involving isoscalar operator insertions, where
disconnected loops with all momentum components of the quark field are needed.
A well-known example of such a computation would be the evaluation of the
strangeness content of the nucleon.
